\section{AHSME 1964}
        \begin{enumerate}
        	\item (\textit{AHSME 1964}) What is the value of $[\log_{10}(5\log_{10}100)]^2$?


 $\textbf{(A)}\ \log_{10}50 \qquad \textbf{(B)}\ 25\qquad \textbf{(C)}\ 10 \qquad \textbf{(D)}\ 2\qquad \textbf{(E)}\ 1$

 \textbf{Answer: }E



\textbf{Solution}

$[\log_{10}(5\log_{10}100)]^2$

 Since $10^2 = 100$, we have $\log_{10} 100 = 2$, so:

 $[\log_{10}(5\cdot 2)]^2$

 Multiply:

 $[\log_{10}(10)]^2$

 Since $10^1 = 10$, we have $\log_{10} 10 = 1$, so:

 $[1]^2$
$1$
$\fbox{E}$.

 


	\item (\textit{AHSME 1964}) The graph of $x^2-4y^2=0$ is:


 $\textbf{(A)}\ \text{a parabola} \qquad \textbf{(B)}\ \text{an ellipse} \qquad \textbf{(C)}\ \text{a pair of straight lines}\qquad \\ \textbf{(D)}\ \text{a point}\qquad \textbf{(E)}\ \text{None of these}$

 \textbf{Answer: }C



\textbf{Solution}

In the equation $x^2 - 4y^2 = k$, because the coefficients of $x^2$ and $y^2$ are of opposite sign, the graph is typically a hyperbola for most real values of $k$.  However, there is one exception.  When $k=0$, the equation can be factored as $(x - 2y)(x+2y) = 0$.  This gives the graph of two lines passing though the origin:  $x=2y$ and $x=-2y$.  Thus, the answer is $\boxed{\textbf{(C)}}$



	\item (\textit{AHSME 1964}) When a positive integer $x$ is divided by a positive integer $y$, the quotient is $u$ and the remainder is $v$, where $u$ and $v$ are integers. 
What is the remainder when $x+2uy$ is divided by $y$?


 $\textbf{(A)}\ 0 \qquad \textbf{(B)}\ 2u \qquad \textbf{(C)}\ 3u \qquad \textbf{(D)}\ v \qquad \textbf{(E)}\ 2v$

 \textbf{Answer: }D



\textbf{Solution 1}


            \begin{itemize}
        \item We can solve this problem by elemetary modular arthmetic,
            \end{itemize}
            $x \equiv v\ (\textrm{mod}\ y)$ $=>$ $x+2uy \equiv v\ (\textrm{mod}\ y)$




\textbf{Solution 2}

By the definition of quotient and remainder, problem states that $x = uy + v$.

 The problem asks to find the remainder of $x + 2uy$ when divided by $y$.  Since $2uy$ is divisible by $y$, adding it to $x$ will not change the remainder.  Therefore, the answer is $\boxed{\textbf{(D)}}$.

 



\textbf{Solution 3}

If the statement is true for all values of $(x, y, u, v)$, then it must be true for a specific set of $(x, y, u, v)$.

 If you let $x=43$ and $y = 8$, then the quotient is $u = 5$ and the remainder is $v = 3$.  The problem asks what the remainder is when you divide $x + 2uy = 43 + 2 \cdot 5 \cdot 8 = 123$ by $8$.  In this case, the remainder is $3$.

 When you plug in $u=5$ and $v = 3$ into the answer choices, they become $0, 5, 10, 3, 6$, respectively.  Therefore, the answer is $\boxed{\textbf{(D)}}$.

 


	\item (\textit{AHSME 1964}) The expression


 
 \[\frac{P+Q}{P-Q}-\frac{P-Q}{P+Q}\]


 where $P=x+y$ and $Q=x-y$, is equivalent to:


 $\textbf{(A)}\ \frac{x^2-y^2}{xy}\qquad \textbf{(B)}\ \frac{x^2-y^2}{2xy}\qquad \textbf{(C)}\ 1 \qquad \textbf{(D)}\ \frac{x^2+y^2}{xy}\qquad \textbf{(E)}\ \frac{x^2+y^2}{2xy}$

 \textbf{Answer: }A



\textbf{Solution 1}

Simplifying before substituting by obtaining a common denominator gives: 

 $\frac{(P+Q)^2 - (P - Q)^2}{(P-Q)(P+Q)}$
$\frac{(P^2 + 2PQ + Q^2) - (P^2 - 2PQ + Q^2)}{P^2 - Q^2}$
$\frac{4PQ}{P^2 - Q^2}$

 Substituting gives:

 $\frac{4(x+y)(x-y)}{(x+y)^2 - (x - y)^2}$
$\frac{4(x^2 - y^2)}{(x^2 + 2xy + y^2) - (x^2 - 2xy + y^2)}$
$\frac{4(x^2 - y^2)}{4xy}$
$\frac{x^2 - y^2}{xy}$
$\boxed{\textbf{(A)}}$




\textbf{Solution 2}

Substituting and then simplifying, noting that $P+Q = 2x$ and $P - Q = 2y$, gives:

 $\frac{2x}{2y} - \frac{2y}{2x}$
$\frac{x^2}{xy} - \frac{y^2}{xy}$
$\frac{x^2 - y^2}{xy}$
$\boxed{\textbf{(A)}}$




\textbf{Solution 3}

If it's true for all values of $x$ and $y$, then it's true for a specific value of $x$ and $y$.  Letting $x=5$ and $y = 4$ gives $P = 9$ and $Q = 1$.  The fraction becomes:

 $\frac{9 + 1}{9 - 1} - \frac{9 - 1}{9 + 1}$
$\frac{10}{8} - \frac{8}{10}$
$\frac{5}{4} - \frac{4}{5}$
$\frac{25}{20} - \frac{16}{20}$
$\frac{9}{20}$

 Plugging in $(x, y) = (5, 4)$ into the answer choices gives $\frac{9}{20}, \frac{9}{40}, 1, \frac{41}{20}, \frac{41}{40}$.  Therefore, the answer is $\boxed{\textbf{(A)}}$



	\item (\textit{AHSME 1964}) If $y$ varies directly as $x$, and if $y=8$ when $x=4$, the value of $y$ when $x=-8$ is:


 $\textbf{(A)}\ -16 \qquad \textbf{(B)}\ -4 \qquad \textbf{(C)}\ -2 \qquad \textbf{(D)}\ 4k, k= \pm1, \pm2, \dots \qquad \\ \textbf{(E)}\ 16k, k=\pm1,\pm2,\dots$

 \textbf{Answer: }A



\textbf{Solution}

If $y$ varies directly as $x$, then $\frac{x}{y}$ will be a constant for every pair of $(x, y)$.  Thus, we have:

 $\frac{4}{8} = \frac{-8}{y}$
$y \cdot \frac{1}{2} = -8$
$y = -16$
$\boxed{\textbf{(A)}}$



	\item (\textit{AHSME 1964}) If $x, 2x+2, 3x+3, \dots$ are in geometric progression, the fourth term is:


 $\textbf{(A)}\ -27 \qquad \textbf{(B)}\ -13\frac{1}{2} \qquad \textbf{(C)}\ 12\qquad \textbf{(D)}\ 13\frac{1}{2}\qquad \textbf{(E)}\ 27$

 \textbf{Answer: }B



\textbf{Solution}

Since we know the sequence is a geometric sequence, the ratio of consecutive terms is always the same number.  Thus, we can set up an equation:

 $\frac{2x+2}{x} = \frac{3x+3}{2x+2}$.

 Solving it, we get:

 $\frac{2x+2}{x} = \frac{3x+3}{2x+2}$
$(2x+2)(2x+2) = (3x+3)(x)$
$4x^2+8x+4 = 3x^2+3x$
$x^2+5x+4 = 0$
$(x + 4)(x+1) = 0$
$x = -4$ or $x = -1$

 If $x=-1$, the sequence has a $0$ as the second term, which is not allowed in a geometric sequence, so it is an extraneous solution that came about because we cross-multiplied by $(2x + 2)$, which is $0$.

 If $x=-4$, we plug into $x, 2x +2, 3x + 3$ to find the sequence starts as $-4, -6, -9$.  The common ratio is $\frac{-6}{-4} = \frac{3}{2}$.  The next term is $\frac{3}{2} \cdot -9 = \frac{-27}{2} = -13\frac{1}{2}$, which is option $\textbf{(B)}$



	\item (\textit{AHSME 1964}) Let n be the number of real values of $p$ for which the roots of 
$x^2-px+p=0$
are equal. Then n equals:


 $\textbf{(A)}\ 0 \qquad \textbf{(B)}\ 1 \qquad \textbf{(C)}\ 2 \qquad \textbf{(D)}\ \text{a finite number greater than 2}\qquad \textbf{(E)}\ \infty$

 \textbf{Answer: }C



\textbf{Solution}

The quadratic $x^2 - px + p$ has discriminant $p^2 - 4p = p(p-4)$.  For the roots of this quadratic to be equal, the discriminant must be $0$, which only happens if $p = 0$ or $p = 4$.  Therefore, there are $\boxed{\textbf{(C)}\ 2}$ such values of $p$.

 


	\item (\textit{AHSME 1964}) The smaller root of the equation 
$\left(x-\frac{3}{4}\right)\left(x-\frac{3}{4}\right)+\left(x-\frac{3}{4}\right)\left(x-\frac{1}{2}\right) =0$ is:


 $\textbf{(A)}\ -\frac{3}{4}\qquad \textbf{(B)}\ \frac{1}{2}\qquad \textbf{(C)}\ \frac{5}{8}\qquad \textbf{(D)}\ \frac{3}{4}\qquad \textbf{(E)}\ 1$

 \textbf{Answer: }C



\textbf{Solution}

Note that this equation is of the form $a^2 + ab = 0$, which factors to $a(a + b) = 0$.  Plugging in $a = x - \frac{3}{4}$ and $b = x - \frac{1}{2}$ gives:

 $(x - \frac{3}{4})(x - \frac{3}{4} + x - \frac{1}{2}) = 0$
$(x - \frac{3}{4})(2x - \frac{5}{4}) = 0$

 The roots are $x = \frac{3}{4}$ nad $x = \frac{1}{2} \cdot \frac{5}{4}$.  The smaller root is $\frac{5}{8}$, which is option $\boxed{\textbf{(C)}}$.

 


	\item (\textit{AHSME 1964}) A jobber buys an article at $\$24$ less $12\frac{1}{2}\%$. He then wishes to sell the article at a gain of $33\frac{1}{3}\%$ of his cost 
after allowing a $20\%$ discount on his marked price. At what price, in dollars, should the article be marked?


 $\textbf{(A)}\ 25.20 \qquad \textbf{(B)}\ 30.00 \qquad \textbf{(C)}\ 33.60 \qquad \textbf{(D)}\ 40.00 \qquad \textbf{(E)}\ \text{none of these}$

 \textbf{Answer: }E



\textbf{Solution}

The item is bought for $24$ minus $12\frac{1}{2}\%$ of $24$.  Since $12\frac{1}{2}\% = \frac{1}{8}$, the item was purchased for $24 - 24 \cdot \frac{1}{8}$, or $21$.

 He wants to make $33\frac{1}{3}\%$ of the purchase price, or $21 \cdot \frac{1}{3} = 7$.  This means the price must actually sell for $21 + 7 = 28$.

 If the article were marked for $x$ dollars, then after the $20\%$ discount, it would sell for $\frac{4}{5}x$.  Thus, we have $\frac{4}{5}x = 28$, which leads to $x = 28 \cdot \frac{5}{4}$, or $35$ dollars.  The correct answer is $\boxed{\textbf{(E)}}$



	\item (\textit{AHSME 1964}) Given a square side of length $s$. On a diagonal as base a triangle with three unequal sides is constructed so that its area
equals that of the square. The length of the altitude drawn to the base is:


 $\textbf{(A)}\ s\sqrt{2} \qquad \textbf{(B)}\ s/\sqrt{2} \qquad \textbf{(C)}\ 2s \qquad \textbf{(D)}\ 2\sqrt{s} \qquad \textbf{(E)}\ 2/\sqrt{s}$

 \textbf{Answer: }A



\textbf{Solution}

The area of the square is $s^2$.  The diagonal of a square with side $s$ bisects the square into two $45-45-90$ right triangles, so the diagonal has length $s\sqrt{2}$.

 The area of the triangle is $\frac{1}{2}bh$.  The base $b$ of the triangle is the diagonal of the square, which is $b = s\sqrt{2}$.  If the area of the triangle is equal to the area of the square, we have:

 $s^2 = \frac{1}{2}bh$
$s^2 = \frac{1}{2}s\sqrt{2}\cdot h$
$s = \frac{\sqrt{2}}{2}h$
$h = \frac{2}{\sqrt{2}}s$
$h = s\sqrt{2}$

 This is option $\boxed{\textbf{(A)}}$



	\item (\textit{AHSME 1964}) Given $2^x=8^{y+1}$ and $9^y=3^{x-9}$, find the value of $x+y$
$\textbf{(A)}\ 18 \qquad \textbf{(B)}\ 21 \qquad \textbf{(C)}\ 24 \qquad \textbf{(D)}\ 27 \qquad \textbf{(E)}\ 30$

 \textbf{Answer: }D



\textbf{Solution}

Since $8^{y + 1} = 2^{3(y+1)}$ and $9^y = 3^{2y}$, we have:

 $2^x = 2^{3(y+1)}$ and $3^{2y} = 3^{x - 9}$

 
Note that if $a^b = a^c$, then $b=c$.  Setting the exponents equal gives $x = 3y + 3$ and $2y = x - 9$.  Plugging the first equation into the second equation gives:

 $2y = (3y + 3) - 9$
$2y = 3y - 6$
$0 = y - 6$
$y = 6$

 Plugging that back in to $x = 3y + 3$ gives $x = 3(6) + 3$, or $x = 21$.  Thus, $x+y = 21 + 6$, or $x+y=27$, which is option $\boxed{\textbf{(D)}}$



	\item (\textit{AHSME 1964}) Which of the following is the negation of the statement: For all $x$ of a certain set, $x^2>0$?


 $\textbf{(A)}\ \text{For all x}, x^2 < 0\qquad \textbf{(B)}\ \text{For all x}, x^2 \le 0\qquad \textbf{(C)}\ \text{For no x}, x^2>0\qquad \\ \textbf{(D)}\ \text{For some x}, x^2>0\qquad \textbf{(E)}\ \text{For some x}, x^2 \le 0$

 \textbf{Answer: }E



\textbf{Solution 1}

The given statement is false if and only if there exists $x$ such that $x^2$ is not greater than $0$ (that is, $x \leq 0$).  Thus the negation of the given statement is $\boxed{\textbf{(E)}\ \text{For some }x\text{, }x^2 \leq 0}$.

 (A), (B), (C) are incorrect as they fail to include cases where $x^2 > 0$ for one value of $x$ and $x^2 < 0$ for another value of $x$.

 (D) is incorrect as it fails to include cases where $x^2 \leq 0$ for all $x$.

 



\textbf{Solution 2}

Recall that the negation of $\forall x P(x)$ is $\exists x(\lnot P(x))$.

 Thus the negation of $\forall x(x^2 > 0)$ is $\exists x(x^2 \leq 0)$, or $\boxed{\textbf{(E)}\ \text{For some }x\text{, }x^2 \leq 0}$.

 


	\item (\textit{AHSME 1964}) A circle is inscribed in a triangle with side lengths $8, 13$, and $17$. Let the segments of the side of length $8$, 
made by a point of tangency, be $r$ and $s$, with $r<s$. What is the ratio $r:s$? 


 $\textbf{(A)}\ 1:3 \qquad \textbf{(B)}\ 2:5 \qquad \textbf{(C)}\ 1:2 \qquad \textbf{(D)}\ 2:3 \qquad \textbf{(E)}\ 3:4$

 \textbf{Answer: }A



\textbf{Solution 1}

Label our triangle $ABC$ where that $AB=17$, $AC=13$, and $BC=8$. Let $L$, $J$, and $K$ be the tangency points of $BC$, $AC$, and $AB$ respectively. Let $AK=x$, which implies $AJ=x$. Thus $KB=BL=17-x$ and $JC=LC=13-x$. 

 Since $BL+LC=(17-x)+(13-x)=8$, $x=11$. Thus $r:s=CL:BL=13-x:17-x=2:6=1:3$, hence our answer is $\fbox{A}$.

 



\textbf{Solution 2}

Let the three sides of $8, 13, 17$ be split into two lengths each, for a total of six segments.  Because the vertices of the triangle form two congruent external tangents to the triangle each, there are really only three unique lengths, and two of these three lengths each make up a side.  Thus, we can make a system of three equations in there variables:  $a+b = 8, b+c=13, c+a=17$.  

 Adding all three equations and dividing by $2$ gives $a+b+c = 19$.  Subtracting each equation from that gives $(a, b, c) = (6,2,11)$

 Thus, $a+b$ makes up the shortest side of $8$.  The ratio is $2:6$, or $1:3$, which is answer $\fbox{A}$.

 


	\item (\textit{AHSME 1964}) A farmer bought $749$ sheep. He sold $700$ of them for the price paid for the $749$ sheep. 
The remaining $49$ sheep were sold at the same price per head as the other $700$. 
Based on the cost, the percent gain on the entire transaction is:


 $\textbf{(A)}\ 6.5 \qquad \textbf{(B)}\ 6.75 \qquad \textbf{(C)}\ 7 \qquad \textbf{(D)}\ 7.5 \qquad \textbf{(E)}\ 8$

 \textbf{Answer: }C



\textbf{Solution}

Let us say each sheep cost $x$ dollars.  The farmer paid $749x$ for the sheep.  He sold $700$ of them for $749x$, so each sheep sold for $\frac{749}{700} = 1.07x$.

 Since every sheep sold for the same price per head, and since every sheep cost $x$ and sold for $1.07x$, there is an increase of $\frac{1.07x - 1x}{1x} = 0.07$, or $7\%$, which is option $\boxed{\textbf{(C)}}$



	\item (\textit{AHSME 1964}) A line through the point $(-a,0)$ cuts from the second quadrant a triangular region with area $T$. The equation of the line is:


 $\textbf{(A)}\ 2Tx+a^2y+2aT=0 \qquad \textbf{(B)}\ 2Tx-a^2y+2aT=0 \qquad \textbf{(C)}\ 2Tx+a^2y-2aT=0 \qquad \\ \textbf{(D)}\ 2Tx-a^2y-2aT=0 \qquad \textbf{(E)}\ \text{none of these}$

 \textbf{Answer: }B



\textbf{Solution}

The right triangle has area $T$ and base $a$, so the height $h$ satisfies $\frac{1}{2}ha = T$.  This means $h = \frac{2T}{a}$.  Becuase the triangle is in the second quadrant, the coordinates are the origin, $(-a, 0)$, and $(0, h)$, which means the third coordinate is $(0, \frac{2T}{a})$.

 So, we want the equation of a line though $(-a, 0)$ and $(0, \frac{2T}{a})$.  The slope is $\frac{\frac{2T}{a} - 0}{0 - (-a)}$, which simplifies to $\frac{2T}{a^2}$.

 The y-intercept is $(0, \frac{2T}{a})$, so the line in slope-intercept form is $y = mx + b$, or:

 
$y = \frac{2T}{a^2}x + \frac{2T}{a}$

 All the solutions have a positive x-coefficient and no fractions, so we clear the fractions by multiplying by $a^2$ and move the $y$ term to the right to get:

 $0 = 2Tx + 2aT - a^2y$

 This is equivalent to option $\boxed{\textbf{(B)}}$



	\item (\textit{AHSME 1964}) Let $f(x)=x^2+3x+2$ and let $S$ be the set of integers $\{0, 1, 2, \dots , 25 \}$. 
The number of members $s$ of $S$ such that $f(s)$ has remainder zero when divided by $6$ is:


 $\textbf{(A)}\ 25\qquad \textbf{(B)}\ 22\qquad \textbf{(C)}\ 21\qquad \textbf{(D)}\ 18 \qquad \textbf{(E)}\ 17$

 \textbf{Answer: }E



\textbf{Solution }

Note that for all polynomials $f(x)$, $f(x + 6) \equiv f(x) \pmod 6$.  

 Proof:If $f(x) = a_nx^n + a_{n-1}x^{n-1} + ... + a_0$, then $f(x+6) = a_n(x+6)^n + a_{n-1}(x+6)^{n-1} +...+ a_0$.  In the second equation, we can use the binomial expansion to expand every term, and then subtract off all terms that have a factor of $6^1$ or higher, since subtracting a multiple of $6$ will not change congruence $\pmod 6$.  This leaves $a_nx^n + a_{n-1}x^{n-1} + ... + a_0$, which is $f(x)$, so $f(x+6) \equiv f(x) \pmod 6$.

 So, we only need to test when $f(x)$ has a remainder of $0$ for $0, 1, 2, 3, 4, 5$.  The set of numbers $6, 7, 8, 9, 10, 11$ will repeat remainders, as will all other sets.  The remainders are $2, 0, 0, 2, 0, 0$.

 This means for $s=1, 2, 4, 5$, $f(s)$ is divisible by $6$.  Since $f(1)$ is divisible, so is $f(s)$ for $s=7, 13, 19, 25$, which is $5$ values of $s$ that work.  Since $f(2)$ is divisible, so is $f(s)$ for $s=8, 14, 20$, which is $4$ more values of $s$ that work.  The values of $s=4, 5$ will also generate $4$ solutions each, just like $f(2)$.  This is a total of $17$ values of $s$, for an answer of $\boxed{\textbf{(E)}}$




\textbf{Solution 2}


            \begin{itemize}
        \item We know that,
            \end{itemize}
            $f(x)$ = $x^2$ +$3x$ + $2$ is $0$ congruent modulo 6 that implies $x+1$ or/and $x+2$ is $0$ congruent modulo $6$.The numbers are of the form $6k+5$ and $6k+4$ for some integer $k$ and due to restrictions of number of elements in the set $S$ we get the inequality $k<4$.Then we calculate the possible combinations to get  $17$ as answer i.e. option $\boxed{\textbf{(E)}}$.

 \begin{verbatim}
																																			$Solution$ $by$ $GEOMETRY-WIZARD$
\end{verbatim}

THE above solution does not give you the answer there are more cases,$f(x)$ = $x^2$ + $3x$ + $2$ is congruent to $0$ modulo 6 this has 4 cases,When, 

 
            \begin{itemize}
        \item $x+1$ \equiv $0$(mod 6)
        \item $x+2$ \equiv $0$(mod 6)
        \item $x+1$ \equiv $2$(mod 6), $x+2$ \equiv $3$ (mod 6)
        \item $x+1$ \equiv $0$ (mod 6), then $x+2$ \equiv $4$ (mod 6), which satisfies.
            \end{itemize}
            Therefore by solvin' these cases we get 17.

 \begin{verbatim}
$EDITED$ $by$ $RNVAA$
\end{verbatim}




	\item (\textit{AHSME 1964}) Given the distinct points $P(x_1, y_1), Q(x_2, y_2)$ and $R(x_1+x_2, y_1+y_2)$. 
Line segments are drawn connecting these points to each other and to the origin $O$. 
Of the three possibilities: (1) parallelogram (2) straight line (3) trapezoid, figure $OPRQ$, 
depending upon the location of the points $P, Q$, and $R$, can be:


 $\textbf{(A)}\ \text{(1) only}\qquad \textbf{(B)}\ \text{(2) only}\qquad \textbf{(C)}\ \text{(3) only}\qquad \textbf{(D)}\ \text{(1) or (2) only}\qquad \textbf{(E)}\ \text{all three}$

 \textbf{Answer: }D



\textbf{Solution}

Using vector addition can help solve this problem quickly.  Note that algebraically, adding $\overrightarrow{OP}$ to $\overrightarrow{OQ}$ will give $\overrightarrow{OR}$.  One method of vector addition is literally known as the "parallelogram rule" - if you are given $\overrightarrow{OP}$ and $\overrightarrow{OQ}$, to find $\overrightarrow{OR}$, you can literally draw a parallelogram, making a line though $P$ parallel to $OQ$, and a line through $Q$ parallel to $OP$.  The intersection of those lines will give the fourth point $R$, and that fourth point will form a parallelogram with $O, P, Q$.

 Thus, $1$ is a possibility.  Case $2$ is also a possibility, if $O, P, Q$ are collinear, then $R$ is also on that line.  

 Since $OP \parallel QR$ and $PQ \parallel RO$, which can be seen from either the prior reasoning or by examining slopes, the figure can never be a trapezoid, which requires exactly one of parallel sides.

 Thus, the answer is $\boxed{\textbf{(D)}}$



	\item (\textit{AHSME 1964}) Let $n$ be the number of pairs of values of $b$ and $c$ such that $3x+by+c=0$ and $cx-2y+12=0$ have the same graph. Then $n$ is:


 $\textbf{(A)}\ 0\qquad \textbf{(B)}\ 1\qquad \textbf{(C)}\ 2\qquad \textbf{(D)}\ \text{finite but more than 2}\qquad \textbf{(E)}\ \infty$

 \textbf{Answer: }C



\textbf{Solution}

For two lines to be the same, their slopes must be equal and their intercepts must be equal.  This is a necessary and sufficient condition.

 The slope of the first line is $\frac{-3}{b}$, while the slope of the second line is $\frac{c}{2}$.  Thus, $\frac{-3}{b} = \frac{c}{2}$, or $bc = -6$.

 The intercept of the first line is $\frac{-c}{b}$, while the intercept of the second line is $6$.  Thus, $6 = \frac{-c}{b}$, or $-6b = c$.

 Plugging $-6b = c$ into $bc = -6$ gives $b(-6b) = -6$, or $b^2 = 1$.  This means $b = \pm 1$  This in turn gives $c = \mp 6$.  Thus, $(b, c) = (\pm 1, \mp 6)$, for two solutions, which is answer $\boxed{\textbf{(C)}}$



	\item (\textit{AHSME 1964}) If $2x-3y-z=0$ and $x+3y-14z=0, z \neq 0$, the numerical value of $\frac{x^2+3xy}{y^2+z^2}$ is:


 $\textbf{(A)}\ 7\qquad \textbf{(B)}\ 2\qquad \textbf{(C)}\ 0\qquad \textbf{(D)}\ -20/17\qquad \textbf{(E)}\ -2$

 \textbf{Answer: }A



\textbf{Solution 1}

If the value of $\frac{x^2+3xy}{y^2+z^2}$ is constant, as the answers imply, we can pick a value of $z$, and then solve the two linear equations for the corresponding $(x, y)$.  We can then plug in $(x,y, z)$ into the expression to get the answer.

 If $z=1$, then $2x - 3y = 1$ and $x + 3y = 14$.  We can solve each equation for $3y$ and set them equal, which leads to $2x - 1 = 14 - x$.  This leads to $x = 5$.  Plugging in $x=5$ into $x + 3y = 14$ gives $y = 3$.  Thus, $(5, 3, 1)$ is one solution to the intersection of the two planes given.

 Plugging $(5, 3, -1)$ into the expression gives $\frac{x^2+3xy}{y^2+z^2}$ gives $\frac{25 + 45}{10}$, or $7$, which is answer $\boxed{\textbf{(A)}}$




\textbf{Solution 2}

If we think of $z$ as a parameter, we get $2x - 3y = z$ and $x + 3y = 14z$.  Adding the equations leads to $3x = 15z$, or $x = 5z$.  Plugging that into $x + 3y = 14z$ gives $5z + 3y = 14z$, or $y = 3z$.  Thus, the intersection of the two planes is given by the parametric line $(5z, 3z, z)$, where $z$ varies along all real numbers.

 We plug this in to $\frac{x^2+3xy}{y^2+z^2}$ to get $\frac{25z^2 + 45z^2}{10z^2}$, or $7$, which is answer $\boxed{\textbf{(A)}}$.

 


	\item (\textit{AHSME 1964}) The sum of the numerical coefficients of all the terms in the expansion of $(x-2y)^{18}$ is:


 $\textbf{(A)}\ 0\qquad \textbf{(B)}\ 1\qquad \textbf{(C)}\ 19\qquad \textbf{(D)}\ -1\qquad \textbf{(E)}\ -19$

 \textbf{Answer: }B



\textbf{Solution}

For any polynomial, even a polynomial with more than one variable, the sum of all the coefficients (including the constant, which is the coefficient of $x^0y^0$) is found by setting all variables equal to $1$.  Note that we don't have to worry about whether a constant is a coefficient of an "invisible" $x^0y^0$ term, because there is no such term here.

 Setting $x=y=1$ gives $(-1)^{18}$, which is equal to $1$, which is answer $\boxed{\textbf{(B)}}$.

 


	\item (\textit{AHSME 1964}) If $\log_{b^2}x+\log_{x^2}b=1, b>0, b \neq 1, x \neq 1$, then $x$ equals:


 $\textbf{(A)}\ 1/b^2 \qquad \textbf{(B)}\ 1/b \qquad \textbf{(C)}\ b^2 \qquad \textbf{(D)}\ b \qquad \textbf{(E)}\ \sqrt{b}$

 \textbf{Answer: }D



\textbf{Solution 1}

Using natural log as a "neutral base", and applying the change of base formula to each term, we get:

 $\frac{\ln x}{\ln b^2} + \frac{\ln b}{\ln x^2} = 1$
$\frac{\ln x}{2\ln b} + \frac{\ln b}{2\ln x} = 1$
$\frac{\ln x \ln x + \ln b \ln b}{2\ln b \ln x} = 1$
$\ln x \ln x + \ln b \ln b = 2\ln b \ln x$

 You could inspect the equation here and see that $x=b$ is one solution.  Or, you can substitute $X = \ln x$ and $B = \ln b$ to get a quadratic in $X$:

 $X^2 + B^2 = 2BX$
$X^2 - 2BX + B^2 = 0$

 The above is a quadratic with coefficients $(1, -2B, B^2)$.  Plug into the QF to get:

 $X = \frac{2B \pm \sqrt{4B^2 - 4B^2}}{2}$
$X = B$
$\ln x = \ln b$
$x = b$

 Either way, the answer is $\boxed{\textbf{(D)}}$.

 



\textbf{Solution 2}

All answers are of the form $x = b^n$, so we substitute that into the equation and try to solve for $n$.  We get:

 $\log_{b^2}x+\log_{x^2}b=1$
$\log_{b^2}b^n + \log_{b^{2n}} b = 1$

 By the definition of a logarithm, the first term on the left is asking for the exponent $x$ needed to change the number $b^2$ into $(b^2)^x$ to get to $b^n$.  That exponent is $\frac{n}{2}$.

 The second term is asking for a similar exponent needed to change $b^{2n}$ into $b$.  That exponent is $\frac{1}{2n}$.

 The equation becomes $\frac{n}{2} + \frac{1}{2n} = 1$.  Multiplying by $2n$ gives the quadratic $n^2 + 1 = 2n$, which has the solution $n=1$.  Thus, $x = b^n = b^1$, and the answer is $\boxed{\textbf{(D)}}$.

 


	\item (\textit{AHSME 1964}) Given parallelogram $ABCD$ with $E$ the midpoint of diagonal $BD$. Point $E$ is connected to a point $F$ in $DA$ so that 
$DF=\frac{1}{3}DA$. What is the ratio of the area of $\triangle DFE$ to the area of quadrilateral $ABEF$?


 $\textbf{(A)}\ 1:2 \qquad \textbf{(B)}\ 1:3 \qquad \textbf{(C)}\ 1:5 \qquad \textbf{(D)}\ 1:6 \qquad \textbf{(E)}\ 1:7$

 \textbf{Answer: }C



\textbf{Solution}

If it works for a parallelogram $ABCD$, it should also work for a unit square, with $A(0, 0), B(0, 1), C(1, 1), D(1, 0)$.  We are given that $E$ is the midpoint of $BD$, so $E(0.5, 0.5)$.  If $F$ is on $DA$, then $F(x, 0)$.  We note that $DF = 1-x$ and $DA = 1$, so $DF = \frac{1}{3}DA$ means $1-x = \frac{1}{3}$, or $x = \frac{2}{3}$, and hence $F(\frac{2}{3}, 0)$.

 We note that $\triangle DFE$ has a base $DF$ that is $\frac{1}{3}$ and an altitude from $E$ to $DF$ that is $\frac{1}{2}$.  Therefore, $[DEF] = \frac{1}{2}bh = \frac{1}{2} \cdot \frac{1}{3} \cdot \frac{1}{2} = \frac{1}{12}$.

 Quadrilateral $ABEF$ can be split into $\triangle ABE$ and $\triangle AEF$.  The first triangle is $\frac{1}{4}$ of the unit square cut diagonally, so $[ABE] = \frac{1}{4}$.  The second triangle has base $AF$ that is $\frac{2}{3}$ and height $E$ to $AF$ that is $\frac{1}{2}$.  Therefore, $[AEF] = \frac{1}{2}bh = \frac{1}{2} \cdot \frac{2}{3} \cdot \frac{1}{2} = \frac{1}{6}$.

 The entire quadrilateral $ABEF$ has area $\frac{1}{4} + \frac{1}{6} = \frac{5}{12}$.  This is $5$ times larger than the area of $\triangle DFE$, so the ratio is $1:5$, or $\boxed{\textbf{(C)}}$.

 



\textbf{Solution 2}


 \begin{align*} \frac{A_{DFE}}{A_{ADE}} &= \frac{\frac{1}{3}DA}{DA} = \frac{1}{3} \\ A_{DFE} &= \frac{1}{3}A_{ADE} \\ A_{ADE} &= \frac{1}{2}A_{ADB} = \frac{1}{4}A_{ABCD} \\ \implies A_{DFE} &= \frac{1}{3} \cdot \frac{1}{4} A_{ABCD} = \frac{1}{12} A_{ABCD} \\ A_{ABEF} &= A_{ABD} - A_{DFE} \\ &= \frac{1}{2}A_{ABCD} - \frac{1}{12}A_{ABCD} \\ &= \frac{5}{12}A_{ABCD} \end{align*}
Therefore, $\frac{A_{DFE}}{A_{ABEF}} = \frac{\frac{1}{12}}{\frac{5}{12}} = \frac{1}{5}$, giving us the answer $\boxed{\textbf{(C)}}$. -nullptr07

 


	\item (\textit{AHSME 1964}) Two numbers are such that their difference, their sum, and their product are to one another as $1:7:24$. The product of the two numbers is:


 $\textbf{(A)}\ 6\qquad \textbf{(B)}\ 12\qquad \textbf{(C)}\ 24\qquad \textbf{(D)}\ 48\qquad \textbf{(E)}\ 96$

 \textbf{Answer: }D



\textbf{Solution}

Set the two numbers as $x$ and $y$. Therefore, $x+y=7(x-y),  xy=24(x-y)$, and$24(x+y)=7xy$. Simplifying the first equation gives $y=\frac{3}{4}x$. Substituting for $y$ in the second equation gives $\frac{3}{4}x^2=6x.$ Solving yields $x=8$ or $x=0$. Substituting $x=0$ back into the first equation yields $1=-7$ which is false, so $x=0$ is not valid and $x=8$. Substituting into $y=\frac{3}{4}x$ gives $y=6$ and $xy=\boxed{\textbf{(D) } 48}$.

 \textbf{See More} 


	\item (\textit{AHSME 1964}) Let $y=(x-a)^2+(x-b)^2, a, b$ constants. For what value of $x$ is $y$ a minimum?


 $\textbf{(A)}\ \frac{a+b}{2} \qquad \textbf{(B)}\ a+b \qquad \textbf{(C)}\ \sqrt{ab} \qquad \textbf{(D)}\ \sqrt{\frac{a^2+b^2}{2}}\qquad \textbf{(E)}\ \frac{a+b}{2ab}$

 \textbf{Answer: }A



\textbf{Solution 1}

Expanding the quadratic and collecting terms gives $y = 2x^2 - (2a + 2b)x + (a^2+b^2)$.  For a quadratic of the form $y = Ax^2 + Bx + C$ with $A>0$, $y$ is minimized when $x = -\frac{B}{2A}$, which is the average of the roots.

 Thus, the quadratic is minimized when $x = \frac{2a+2b}{2\cdot 2} = \frac{a+b}{2}$, which is answer $\boxed{\textbf{(A)}}$.

 



\textbf{Solution 2}

The problem should return real values for $ab = 0$ and $ab < 0$, which eliminates $E$ and $C$.  We want to distinguish between options $A, B, D$, and testing $(a, b) = (2, 0)$ should do that, as answers $A, B, D$ will turn into $1, 2, \sqrt{2}$, respectively.

 PLugging in $(a, b) = (2, 0)$ gives $y = (x - 2)^2 + x^2$, or $y = 2x^2 - 4x + 4$.  This has a minimum at $x = -\frac{-4}{4}$, or at $x=1$.  This is answer $\boxed{\textbf{(A)}}$.

 


	\item (\textit{AHSME 1964}) The set of values of $m$ for which $x^2+3xy+x+my-m$ has two factors, with integer coefficients, which are linear in $x$ and $y$, is precisely:


 $\textbf{(A)}\ 0, 12, -12\qquad \textbf{(B)}\ 0, 12\qquad \textbf{(C)}\ 12, -12\qquad \textbf{(D)}\ 12\qquad \textbf{(E)}\ 0$

 \textbf{Answer: }B



\textbf{Solution}

Since there are only $3$ candidate values for $m$, we test $0, 12, -12$.

 If $m=0$, then the expression is $x^2 + 3xy + x$.  The term $x$ appears in each monomial, giving $x(x + 3y + 1)$, which has two factors that are linear in $x$ and $y$ with integer coefficients.  This eliminates $C$ and $D$.

 We next test $m=12$.  This gives $x^2 + 3xy + x + 12y - 12$.  The linear factors will be of the form $(Ax + By + C)(Dx + Ey + F)$.  We will match coefficients from the skeleton form to the expression that is to be factored.

 Matching the $y^2$ coefficient gives $BE = 0$, and WLOG we can let $B=0$ to give $(Ax + C)(Dx + Ey + F)$.  

 Matching the $x^2$ coefficients gives $AD = 1$.  Since $A, D$ are integers, we either have $(A, D) = (1, 1), (-1, -1)$.  WLOG we can pick the former, since if $(a)(b)$ is a factorization, so is $(-a)(-b)$.  

 We now have $(x + C)(x + Ey + F)$.  Matching the $xy$ term gives $E =3$.  We now have $(x+C)(x + 3y + F)$.  Matching the $y$ term gives $3C = 12$, which leads to $C = 4$.  Finally, matching constants leads to $4F = -12$ and $F = -3$ for a provisional factorization of $(x + 4)(x + 3y - 3)$.  We've matched every coefficient except for $x$, and finding the $x$ term by selectively multiplying leads to $4x + -3x = x$, which matches, so this is a real factorization.

 We now must check $m=-12$.  The expression is $x^2+3xy+x-12y+12m$, and the skeleton once again is $(Ax + By + C)(Dx + Ey + F)$.  The beginning three steps matching $y^2, x^2, xy$ are the same, leading to $(x + C)(x + 3y + F)$.  Matching the $y$ term gives $3C = -12$, or $C = -4$.  The skeleton becomes $(x -4)(x + 3y + F)$.  The constant locks in $F = -3$ for a provisional factorization of $(x-4)(x + 3y - 3)$.  However, the $x$ term does not match up, because it is $-4x + -3x = -7x$, when it needs to be $x$.

 Thus, $m=0, 12$ work and $m=-12$ does not, giving answer $\boxed{\textbf{(B)}}$.

 


	\item (\textit{AHSME 1964}) In a ten-mile race $\textit{First}$ beats $\textit{Second}$ by $2$ miles and 
$\textit{First}$ beats $\textit{Third}$ by $4$ miles. 
If the runners maintain constant speeds throughout the race, 
by how many miles does $\textit{Second}$ beat $\textit{Third}$?


 $\textbf{(A)}\ 2\qquad \textbf{(B)}\ 2\frac{1}{4}\qquad \textbf{(C)}\ 2\frac{1}{2}\qquad \textbf{(D)}\ 2\frac{3}{4}\qquad \textbf{(E)}\ 3$

 \textbf{Answer: }C



\textbf{Solution}

At the time First finishes the race, Second has run $10-2=8$ miles while Third has run $10-4=6$ miles.  Thus Third's speed is $\frac{6}{8} = \frac{3}{4}$ of Second's speed.

 Therefore, when Second finishes the race, Third has run $\frac{3}{4}\cdot 10 = \frac{15}{2}$ miles.  So Second beats Third by $10 - \frac{15}{2} = \frac{5}{2} = \boxed{\textbf{(C)}\ 2\frac{1}{2}}$ miles.

 


	\item (\textit{AHSME 1964}) If $x$ is a real number and $|x-4|+|x-3|<a$ where $a>0$, then:


 $\textbf{(A)}\ 0<a<.01\qquad \textbf{(B)}\ .01<a<1 \qquad \textbf{(C)}\ 0<a<1\qquad \\ \textbf{(D)}\ 0<a \le 1\qquad \textbf{(E)}\ a>1$

 \textbf{Answer: }E



\textbf{Solution 1}

Note that $|x-4| = |-(x-4)| = |4-x|$.  

 By the triangle inequality, $|x-4| + |x-3| = |4-x| + |x-3| \geq |(4-x)+(x-3)| = 1$.  Therefore, $\boxed{\textbf{(E)}\ a>1}$.

 



\textbf{Solution 2}

Let $A$, $B$, and $X$ be at $3$, $4$, and $x$ on a number line.  Then $|x-4| + |x-3| = AX + BX$.  

 If $3 \leq x \leq 4$, then $X \in \overline{AB}$, so $AX + BX = AB = 1$.  

 If $x < 3$, then $X$ is to the left of $A$, so $AX + BX = AX + (AX + AB) = 2AX + 1 > 1$.  

 If $x > 4$, then $X$ is to the right of $B$, so $AX + BX = (AB + BX) + BX = 2BX + 1 > 1$.

 So $1 \leq AX + BX < a$ for all $x$, and thus $\boxed{\textbf{(E)}\ a > 1}$ is the correct choice.

 



\textbf{Solution 3 (Answer Choices)}

If $x = 5$, then $a > |x-4|+|x-3|=3$.  Of the choices, only $\boxed{\textbf{(E)}\ a>1}$ is satisfied by $a=3$.

 


	\item (\textit{AHSME 1964}) The sum of $n$ terms of an arithmetic progression is $153$, and the common difference is $2$. 
If the first term is an integer, and $n>1$, then the number of possible values for $n$ is:


 $\textbf{(A)}\ 2\qquad \textbf{(B)}\ 3\qquad \textbf{(C)}\ 4\qquad \textbf{(D)}\ 5\qquad \textbf{(E)}\ 6$

 \textbf{Answer: }D



\textbf{Solution}

Let the progression start at $a$, have common difference $2$, and end at $a + 2(n-1)$.

 
The average term is $\frac{a + (a + 2(n-1))}{2}$, or $a + n - 1$.  Since the number of terms is $n$, and the sum of the terms is $153$, we have:

 $n(a+n-1) = 153$

 Since $n$ is a positive integer, it must be a factor of $153$.  This means $n = 1, 3, 9, 17, 51, 153$ are the only possibilities.  We are given $n>1$, leaving the other five factors.

 We now must check if $a$ is an integer.  We have $a = \frac{153}{n} + 1 - n$.  If $n$ is a factor of $153$, then $\frac{153}{n}$ will be an integer.  Adding $1-n$ wil keep it an integer.

 Thus, there are $5$ possible values for $n$, which is answer $\boxed{\textbf{(D)}}$.

 


	\item (\textit{AHSME 1964}) In this figure $\angle RFS = \angle FDR, FD = 4$ inches, $DR = 6$ inches, $FR = 5$ inches, $FS = 7\frac{1}{2}$ inches. 
The length of $RS$, in inches, is:


 
                    \begin{figure}[H]
                        \centering
                        \includegraphics[width=0.5\textwidth]{Problem/AHSME_I_1964_29_Problem_Figure_1.png}
                    \end{figure}
                    
$\textbf{(A)}\ \text{undetermined} \qquad \textbf{(B)}\ 4\qquad \textbf{(C)}\ 5\frac{1}{2} \qquad \textbf{(D)}\ 6\qquad \textbf{(E)}\ 6\frac{1}{4}$

 \textbf{Answer: }E



\textbf{Solution 1}

Note that $\frac{SF}{RD} = \frac{FR}{DF} = \frac{5}{4}$.  Since $\angle SFR = \angle RDF$, by SAS similarity $\triangle SFR \sim \triangle RDF$.  

 Therefore, $\frac{RS}{FR} = \frac{5}{4}$, so $RS = \frac{5}{4}FR = \frac{5}{4} \cdot 5 = \frac{25}{4} = \boxed{ \textbf{(E)}  \ 6\frac{1}{4}}$.

 -edited by j314andrews

 



\textbf{Solution 2}

Let $\angle RFS = \angle FDR = \theta$. By Cosine Law, we have:
 \begin{align*} RF^2 &= DR^2 + DF^2 - 2(DR)(DF)cos(\theta) \\ 5^2 &= 6^2 + 4^2 - 2(6)(4)cos(\theta) \\ cos(\theta) &= \frac{6^2+4^2-5^2}{2(6)(4)} = \frac{9}{16} \end{align*}
Applying Cosine Law again in $\triangle RFS$, we have:
 \begin{align*} RS^2 &= RF^2 + FS^2 - 2(RF)(FS)cos(\theta) \\ RS^2 &= 25 + (\frac{15}{2})^2 - 2(5)(\frac{15}{2})(\frac{9}{16}) \\ RS^2 &= \frac{625}{16} \end{align*}
giving us $RS = \frac{25}{4}$, which is the answer $\boxed{\textbf{(E)}}$. -nullptr07

 


	\item (\textit{AHSME 1964}) If $(7+4\sqrt{3})x^2+(2+\sqrt{3})x-2=0$, the larger root minus the smaller root is:


 $\textbf{(A)}\ -2+3\sqrt{3}\qquad \textbf{(B)}\ 2-\sqrt{3}\qquad \textbf{(C)}\ 6+3\sqrt{3}\qquad \textbf{(D)}\ 6-3\sqrt{3}\qquad \textbf{(E)}\ 3\sqrt{3}+2$

 \textbf{Answer: }D



\textbf{Solution 1}

Dividing the quadratic by $7 + 4\sqrt{3}$ to obtain a monic polynomial will give a linear coefficient of $\frac{2 + \sqrt{3}}{7 + 4\sqrt{3}}$.  Rationalizing the denominator gives:

 $\frac{(2 + \sqrt{3})(7 - 4\sqrt{3})}{7^2 - 4^2 \cdot 3}$
$=\frac{14 - 12 - \sqrt{3}}{49-48}$
$=2 - \sqrt{3}$

 Dividing the constant term by $7 + 4\sqrt{3}$ (and using the same radical conjugate as above) gives:

 $\frac{-2}{7 + 4\sqrt{3}}$
$=-2(7 - 4\sqrt{3})$
$=8\sqrt{3} - 14$

 So, dividing the original quadratic by the coefficient of $x^2$ gives $x^2 + (2 - \sqrt{3})x + 8\sqrt{3} - 14 = 0$

 From the quadratic formula, the positive difference of the roots is $\frac{\sqrt{b^2 - 4ac}}{a}$.  Plugging in gives:

 $\sqrt{(2 - \sqrt{3})^2 - 4(8\sqrt{3} - 14)(1)}$
$=\sqrt{7 - 4\sqrt{3} - 32\sqrt{3} + 56}$
$=\sqrt{63 - 36\sqrt{3}}$
$=3\sqrt{7 - 4\sqrt{3}}$

 Note that if we take $\frac{1}{3}$ of one of the answer choices and square it, we should get $7 - 4\sqrt{3}$.The only answers that are (sort of) divisible by $3$ are $6 \pm 3\sqrt{3}$, so those would make a good first guess.  And given that there is a negative sign underneath the radical, $6 - 3\sqrt{3}$ is the most logical place to start.

 Since $\frac{1}{3}$ of the answer is $2 - \sqrt{3}$, and $(2 - \sqrt{3})^2 = 7 - 4\sqrt{3}$, the answer is indeed $\boxed{\textbf{(D)}}$.

 



\textbf{Solution 2}

\textbf{Submitted by BinouTheGuineaPig} | \textit{A step-by-step solution}

 The original equation can be manipulated as follows.

 $(4+4\sqrt{3}+3)x^2+(2+\sqrt{3})x-2=0$
$(2+\sqrt{3})^2x^2+(2+\sqrt{3})x-2=0$

 Substituting $u = (2+\sqrt{3})x$,

 $\quad u^2+u-2=0$
$\quad (u-1)(u+2)=0$
$\quad u=1$ or $u=-2$

 First root of $x$:

 $\quad (2+\sqrt{3})x_1=1$
$\quad x_1=\frac{1}{2+\sqrt{3}}$
$\quad x_1=2-\sqrt{3}$

 Second root of $x$:

 $\quad (2+\sqrt{3})x_2=-2$
$\quad x_2=\frac{-2}{2+\sqrt{3}}$
$\quad x_2=-2(2-\sqrt{3})$
$\quad x_2=-4+2\sqrt{3}$

 Now, to find which root of $x$ is larger:

 Assume that

 $\quad 2-\sqrt{3}>-4+2\sqrt{3}$, and so

 $\quad 6>3\sqrt{3}$
$\quad 2>\sqrt{3}$
$\quad \sqrt{4}>\sqrt{3}$ which is true. Hence, $x_1>x_2$.

 Finally, finding the difference between the larger and smaller roots of $x$:

 $\quad x_1-x_2$
$\quad =(2-\sqrt{3})-(-4+2\sqrt{3})$
$\quad =6-3\sqrt{3}$

 Therefore, the answer is $\boxed{\textbf{(D)}}$.

 


	\item (\textit{AHSME 1964}) Let $f(n)=\dfrac{5+3\sqrt{5}}{10}\left(\dfrac{1+\sqrt{5}}{2}\right)^n+\dfrac{5-3\sqrt{5}}{10}\left(\dfrac{1-\sqrt{5}}{2}\right)^n$. 
Then $f(n+1)-f(n-1)$, expressed in terms of $f(n)$, equals:


 $\textbf{(A)}\ \dfrac{1}{2}f(n) \qquad \textbf{(B)}\ f(n)\qquad \textbf{(C)}\ 2f(n)+1 \qquad \textbf{(D)}\ f^2(n)\qquad \textbf{(E)}\ \frac{1}{2}(f^2(n)-1)$

 \textbf{Answer: }B



\textbf{Solution}

We compute $f(n+1)$ and $f(n-1)$, while pulling one copy of the exponential part outside:

 $f(n+1) = \dfrac{5+3\sqrt{5}}{10}\cdot \frac{1 + \sqrt{5}}{2}\cdot\left(\dfrac{1+\sqrt{5}}{2}\right)^n+\dfrac{5-3\sqrt{5}}{10}\cdot \frac{1 - \sqrt{5}}{2}\cdot\left(\dfrac{1-\sqrt{5}}{2}\right)^n$
$f(n+1) = \dfrac{20+8\sqrt{5}}{20}\left(\dfrac{1+\sqrt{5}}{2}\right)^n+\dfrac{20-8\sqrt{5}}{20}\left(\dfrac{1-\sqrt{5}}{2}\right)^n$
$f(n-1) = \dfrac{5+3\sqrt{5}}{10}\cdot \frac{2}{1 + \sqrt{5}}\cdot\left(\dfrac{1+\sqrt{5}}{2}\right)^n+\dfrac{5-3\sqrt{5}}{10}\cdot \frac{2}{1 - \sqrt{5}}\cdot\left(\dfrac{1-\sqrt{5}}{2}\right)^n$
$f(n-1) = \dfrac{(5+3\sqrt{5})(1 - \sqrt{5})}{5(1 + \sqrt{5})(1 - \sqrt{5})}\left(\dfrac{1+\sqrt{5}}{2}\right)^n+\dfrac{(5-3\sqrt{5})(1 + \sqrt{5})}{5(1 - \sqrt{5})(1 + \sqrt{5})}\left(\dfrac{1-\sqrt{5}}{2}\right)^n$
$f(n-1) = \dfrac{-10-2\sqrt{5}}{5(-4)}\left(\dfrac{1+\sqrt{5}}{2}\right)^n+\dfrac{-10 + 2\sqrt{5}}{5(-4)}\left(\dfrac{1-\sqrt{5}}{2}\right)^n$
$f(n-1) = \dfrac{10+2\sqrt{5}}{20}\left(\dfrac{1+\sqrt{5}}{2}\right)^n+\dfrac{10 - 2\sqrt{5}}{20}\left(\dfrac{1-\sqrt{5}}{2}\right)^n$

 
Computing $f(n+1) - f(n-1)$ gives:

 $f(n+1) - f(n-1) = \dfrac{10+6\sqrt{5}}{20}\left(\dfrac{1+\sqrt{5}}{2}\right)^n+\dfrac{10-6\sqrt{5}}{20}\left(\dfrac{1-\sqrt{5}}{2}\right)^n$
$f(n+1) - f(n-1) = \dfrac{5+3\sqrt{5}}{10}\left(\dfrac{1+\sqrt{5}}{2}\right)^n+\dfrac{5-3\sqrt{5}}{10}\left(\dfrac{1-\sqrt{5}}{2}\right)^n$
$f(n+1) - f(n-1) = f(n)$

 Thus, the answer is $\boxed{\textbf{(B)}}$.

 



\textbf{Solution 2}

Notice that $\left(\frac{1+\sqrt{5}}{2}\right)^n$ and $\left(\frac{1-\sqrt{5}}{2}\right)^n$ are the characteristics roots for the recurrence relation $F_n = F_{n-1} + F_{n-2}$ (think about Binet's formula). And $f(n)$ is the solution (i.e. $a_n$) to the recurrence relation with constants $a = \frac{5+3\sqrt{5}}{10}$ and $b = \frac{5-3\sqrt{5}}{10}$. Thus, $f(n+1) - f(n-1) = f(n)$, and the answer is $\boxed{\textbf{(B)}}$. -nullptr07

 


	\item (\textit{AHSME 1964}) If $\dfrac{a+b}{b+c}=\dfrac{c+d}{d+a}$, then:


 $\textbf{(A)}\ a\text{ must equal }c \qquad \textbf{(B)}\ a+b+c+d\text{ must equal zero }\qquad \textbf{(C)}\ \text{either }a=c\text{ or }a+b+c+d=0,\text{ or both} \qquad \textbf{(D)}\ a+b+c+d\neq 0\text{ if }a=c \qquad \\ \textbf{(E)}\ a(b+c+d)=c(a+b+d)$

 \textbf{Answer: }C



\textbf{Solution}

Cross-multiplying gives:

 $(a+b)(a+d) = (b+c)(c+d)$
$a^2 + ad + ab + bd = bc + bd + c^2 + cd$
$a^2 + ad + ab   - bc - c^2 - cd = 0$
$a(a + b + d) - c(b+c+d)= 0$

 This looks close to turning into option C, but we don't have a $c$ term in the first parentheses, and we don't have an $a$ term in the second parentheses to allow us to complete the factorization.  However, if we both add $ac$ and subtract $ac$ on the LHS, we get:

 $a(a + b + d) + ac - c(b+c+d) - ca= 0$
$a(a+b+d +c) - c(b+c+d+a) = 0$
$(a-c)(a+b+c+d) = 0$

 This is equivalent to $\boxed{\textbf{(C)}}$



	\item (\textit{AHSME 1964}) $P$ is a point interior to rectangle $ABCD$ and such that $PA=3$ inches, $PD=4$ inches, and $PC=5$ inches. Then $PB$, in inches, equals:


 $\textbf{(A) }2\sqrt{3}\qquad \textbf{(B) }3\sqrt{2}\qquad \textbf{(C) }3\sqrt{3}\qquad \textbf{(D) }4\sqrt{2}\qquad \textbf{(E) }2$

                    \begin{figure}[H]
                        \centering
                        \includegraphics[width=0.5\textwidth]{Problem/AHSME_I_1964_33_Problem_Figure_1.png}
                    \end{figure}
                    

 \textbf{Answer: }B



\textbf{Solution}

From point $P$, create perpendiculars to all four sides, labeling them $a, b, c, d$ starting from going north and continuing clockwise.  Label the length $PB$ as $x$:

 
                    \begin{figure}[H]
                        \centering
                        \includegraphics[width=0.5\textwidth]{Solution/AHSME_I_1964_33_Solution_Figure_1.png}
                    \end{figure}
                    

 We have $a^2 + b^2 = 5^2$ and $c^2 + d^2 = 3^2$, leading to $a^2 + b^2 + c^2 + d^2 = 34$.

 We also have $a^2 + d^2 = 4^2$ and $b^2 + c^2 = x^2$, leading to $a^2 + b^2 + c^2 + d^2 = 16 + x^2$.

 Thus, $34 = 16 + x^2$, or $x = \sqrt{18} = 3\sqrt{2}$, which is option $\boxed{\textbf{(B)}}$



	\item (\textit{AHSME 1964}) If $n$ is a multiple of $4$, the sum $s=1+2i+3i^2+ \ldots +(n+1)i^{n}$, where $i=\sqrt{-1}$, equals:


 $\textbf{(A)}\ 1+i\qquad \textbf{(B)}\ \frac{1}{2}(n+2) \qquad \textbf{(C)}\ \frac{1}{2}(n+2-ni) \qquad \\ \textbf{(D)}\ \frac{1}{2}[(n+1)(1-i)+2]\qquad \textbf{(E)}\ \frac{1}{8}(n^2+8-4ni)$

 \textbf{Answer: }C



\textbf{Solution}

The real part is $1 + 3i^2 + 5i^4 + ...$, which is $1 - 3 + 5 - 7 + ...$.  If $n$ is a multiple of $4$, then we have an odd number of terms in total:  we start with $1$ at $n=0$, then add two more terms to get $1 - 3 + 5$ at $n=4$, etc.  With each successive addition, we're really adding a total of $2$, since $-3 + 5 = 2$, and $-7 + 9 = 2$, etc.  

 At $n = 0$, the sum is $1$, and at $n=4$, the sum is $3$.  Since the sum increases linearly, the real part of the sum is $1 + \frac{1}{2}n$.

 The imaginary part is $\frac{1}{i}(2i + 4i^3 + 6i^5 + ...)$, which is $2 - 4 + 6 - 8 + ...$.  This time, we have an even number of terms.  We group pairs of terms to get $(2-4) + (6-8) + ...$, and notice that each pair gives $-2$.  Again, with $n=0$ the imaginary part is $0$, while with $n=4$ the imaginary part is $-2$.  Since again the sum increases linearly, this means the imaginary part is $-\frac{n}{2}$.

 Combining the real and imaginary parts gives $\frac{n}{2} + 1 - \frac{n}{2}i$, which is equivalent to option $\boxed{\textbf{(C)}}$.

 


	\item (\textit{AHSME 1964}) The sides of a triangle are of lengths $13, 14,$ and $15$. The altitudes of the triangle meet at point $H$. 
If $AD$ is the altitude to the side length $14$, what is the ratio $HD:HA$?     


 
 \textbf{Answer: }B



\textbf{Solution 1}

Using Law of Cosinesand the fact that the ratio equals $\frac{\cos(a)}{\cos(b)\cos(c)}$ you get $\fbox{B) 5:11}$




\textbf{Solution 2 (coordinates)}


                    \begin{figure}[H]
                        \centering
                        \includegraphics[width=0.5\textwidth]{Solution/AHSME_I_1964_35_Solution_Figure_1.png}
                    \end{figure}
                    The reason why we have $HD$ shorter than $HA$ is that all the ratios' left hand side ($HD$) is less than the ratios' right hand side ($HA$).

 We label point $A$ as the origin and point $B$, logically, as $(15,0)$. By Heron's Formula, the area of this triangle is $84.$ Thus the height perpendicular to $AB$ has a length of $11.2,$ and by the Pythagorean Theorem, $AE$ and $EB$ have lengths $6.6$ and $8.4,$ respectively. These lengths tell us that $C$ is at $(6.6,11.2)$.

 The slope of $BC$ is $\dfrac{0-11.2}{15-6.6}=-\dfrac{4}{3},$ and the slope of $AD$ is $\dfrac{3}{4}$ by taking the negative reciprocal of $-\dfrac{4}{3}.$ Therefore, the equation of line $AD$ can best be represented by $y=\dfrac{3}{4}x.$

 We next find the intersection of $CE$ and $AD$. We automatically know the $x$-value; it is just $6.6$ because $CE$ is a straight line hitting $(6.6,0).$ Therefore, the $y$-value is at $\dfrac{3}{4}\times 6.6=4.95.$ Therefore, the intersection between $CE$ and $AD$ is at $(6.6,4.95)$.

 We also need to find the intersection between $BC$ and $AD$. To do that, we know that the line of $AD$ is represented as $y=\dfrac{3}{4}x,$ and the slope of line $BC$ is $-\dfrac{4}{3}.$ We just need to find line $BC$'s y-intercept. So far, we have $y=-\dfrac{4}{3}x+b,$ where $b$ is a real y-intercept. We know that $B$ is located at $(15,0),$ so we plug that into the equation and yield $b=20.$ Therefore, the intersection between the two lines is
 \[\dfrac{3}{4}x=-\dfrac{4}{3}x+20, 9x=-16x+240, 25x=240, x=9.6, y=\dfrac{3}{4}\times 9.6, y=7.2.\]After that, we use the distance formula: $HA$ has a length of 
 \[\sqrt{(6.6-0)^2+(4.95-0)^2}=\sqrt{\dfrac{1089}{25}+\dfrac{9801}{400}}=\sqrt{\dfrac{1089*16+9801}{400}}=\sqrt{\dfrac{27225}{400}}=\sqrt{\dfrac{1089}{16}}=\dfrac{33}{4}=8.25,\] and $HD$ has a length of 
 \[\sqrt{(9.6-6.6)^2+(7.2-4.95)^2}=\sqrt{3^2+(\dfrac{36}{5}-\dfrac{99}{20})^2}=\sqrt{9+\dfrac{81}{16}}=\sqrt{\dfrac{225}{16}}=3.75.\]Thus, we have that $\dfrac{3.75}{8.25}=\dfrac{\frac{15}{4}}{\frac{33}{4}}=\dfrac{15}{33}=\dfrac{5}{11}=\boxed{\bold{B}}.$-OreoChocolate

 



\textbf{Solution 3 (Ceva's Theorem)}


                    \begin{figure}[H]
                        \centering
                        \includegraphics[width=0.5\textwidth]{Solution/AHSME_I_1964_35_Solution_Figure_2.png}
                    \end{figure}
                    

 A consequence of Ceva's Theorem that is sometimes attributed to van Aubel states that:

 
 \[\dfrac{AH}{HD} = \dfrac{AE}{EC} + \dfrac{AF}{FB}\]

 We must then find $AE$ and $AF$.  To find $AF$ note that $\triangle AFC$ and $\triangle BFC$ are both right triangles sharing a common height, $FC$.  Thus

 $AF^2+FC^2=AC^2, \text{and } BF^2+FC^2=BC^2 \implies$$AF^2+FC^2=13^2, \text{and } (15-AF)^2+FC^2=14^2$

 Subtracting the two equations to eliminate the common height term ($FC^2$):

 $(15-AF)^2-AF^2=27 \implies AF=\dfrac{33}{5}$

 A similar computation using $\triangle AEB$ and $\triangle CEB$ gives us:

 $AE^2+EB^2=AB^2, \text{and } CE^2+EB^2=BC^2 \implies$$AE^2+EB^2=15^2, \text{and } (13-AE)^2+EB^2=14^2$
$AE^2-(13-AE)^2=29 \implies AE=\dfrac{99}{13}$

 Returning to our original van Aubel equation yields:

 $\dfrac{AH}{HD} = \dfrac{AE}{EC} + \dfrac{AF}{FB} = \dfrac{\dfrac{99}{13}}{13-\dfrac{99}{13}} + \dfrac{\dfrac{33}{5}}{15-\dfrac{33}{5}}=\dfrac{99}{70}+\dfrac{55}{70}=\dfrac{11}{5}$

 Therefore $HD:HA = \boxed{\textbf{(B) }5:11}$




\textbf{Solution 4 (Similar Triangles)}

Because this triangle is a nice triangle, we can easily find that $AD = 12$, $CD = 5$, and $DB = 9$.From this, we know that:$\triangle BDA$ is a 3:4:5 triangle, and it is similar to $\triangle BEC$ which is similar to $\triangle CDH$.Therefore $HD = CD \cdot 3/4 = 15/4$, and $AH = 12-15/4 = 33/4$,$HD/AH = \dfrac{15/4}{33/4} = \dfrac{5}{11}=\boxed{\bold{B}}$. -ELIZABETHWZYING

 


	\item (\textit{AHSME 1964}) In this figure the radius of the circle is equal to the altitude of the equilateral triangle $ABC$. 
The circle is made to roll along the side $AB$, remaining tangent to it at a variable point $T$ 
and intersecting lines $AC$ and $BC$ in variable points $M$ and $N$, respectively. 
Let $n$ be the number of degrees in arc $MTN$. Then $n$, for all permissible positions of the circle:


 $\textbf{(A) }\text{varies from }30^{\circ}\text{ to }90^{\circ}\quad \textbf{(B) }\text{varies from }30^{\circ}\text{ to }60^{\circ} \\ \textbf{(C) }\text{varies from }60^{\circ}\text{ to }90^{\circ} \quad \textbf{(D) }\text{remains constant at }30^{\circ}\quad \textbf{(E) }\text{remains constant at }60^{\circ}$

                    \begin{figure}[H]
                        \centering
                        \includegraphics[width=0.5\textwidth]{Problem/AHSME_I_1964_36_Problem_Figure_1.png}
                    \end{figure}
                    

 \textbf{Answer: }E



\textbf{Solution}


                    \begin{figure}[H]
                        \centering
                        \includegraphics[width=0.5\textwidth]{Solution/AHSME_I_1964_36_Solution_Figure_1.png}
                    \end{figure}
                    

 Reflect $\triangle ABC$ across a line through $C$ parallel to line $AB$.  Let $A'$, $B'$, $M'$, $N'$, and $T'$ be the reflections of $A$, $B$, $M$, $N$, and $T$ across this line respectively.  By symmetry, $A'B'$ is tangent to the circle at $T'$, and arc $M'T'N'$ also has measure $n^\circ$.  

 By the intersecting chords angle theorem, $\angle MCN = \frac{1}{2}(n^\circ + n^\circ) = n^\circ$.  Since $\triangle ABC$ is equilateral, $n = 60$.  Thus $n$ $\boxed{\textbf{(E)}\ \text{remains constant at }60^{\circ}}$



	\item (\textit{AHSME 1964}) Given two positive number $a, b$ such that $a<b$ , let $A.M.$ be their arithmetic mean and let $G.M.$ be their positive geometric mean. 
Then $A.M.$ minus $G.M.$ is always less than:


 $\textbf{(A) }\dfrac{(b+a)^2}{ab}\qquad \textbf{(B) }\dfrac{(b+a)^2}{8b}\qquad \textbf{(C) }\dfrac{(b-a)^2}{ab} \textbf{(D) }\dfrac{(b-a)^2}{8a}\qquad  \textbf{(E) }\dfrac{(b-a)^2}{8b}$

 \textbf{Answer: }D



\textbf{Solution}

$\boxed{ \textbf{(D)} }$



	\item (\textit{AHSME 1964}) The sides $PQ$ and $PR$ of $\triangle PQR$ are respectively of lengths $4$ inches, and $7$ inches. The median $PM$ is $3\frac{1}{2}$ inches. 
Then $QR$, in inches, is:


 $\textbf{(A) }6\qquad \textbf{(B) }7\qquad \textbf{(C) }8\qquad \textbf{(D) }9\qquad  \textbf{(E) }10$

 \textbf{Answer: }D



\textbf{Solution}

By the \href{https://artofproblemsolving.com/wiki/index.php?title=Median\_Formula&action=edit&redlink=1}{Median Formula}, $PM = \frac12\sqrt{2PQ^2+2PR^2-QR^2}$

 Plugging in the numbers given in the problem, we get
 \[\frac72=\frac12\sqrt{2\cdot4^2+2\cdot7^2-QR^2}\]

 Solving,
 \[7=\sqrt{2(16)+2(49)-QR^2}\]
 \[49=32+98-QR^2\]
 \[QR^2=81\]
 \[QR=9=\boxed{D}\]

 -AOPS81619

 


	\item (\textit{AHSME 1964}) The magnitudes of the sides of $\triangle ABC$ are $a, b$, and $c$, as shown, with $c\le b\le a$. 
Through interior point $P$ and the vertices $A, B, C$, lines are drawn meeting the opposite sides in $A', B', C'$, respectively. 
Let $s=AA'+BB'+CC'$. Then, for all positions of point $P$, $s$ is less than:


 $\textbf{(A) }2a+b\qquad \textbf{(B) }2a+c\qquad \textbf{(C) }2b+c\qquad \textbf{(D) }a+2b\qquad  \textbf{(E) } a+b+c$

                    \begin{figure}[H]
                        \centering
                        \includegraphics[width=0.5\textwidth]{Problem/AHSME_I_1964_39_Problem_Figure_1.png}
                    \end{figure}
                    

 \textbf{Answer: }A



\textbf{Solution 1}

Recall that in any triangle, longer sides oppose larger angles.  Therefore, since $c \leq b \leq a$, $\angle ACB \leq \angle ABC \leq \angle BAC$.

 By the exterior angle theorem, $\angle AA'C > \angle ABC \geq \angle ACB$, so $AA' < AC = b$.

 Similarly, $\angle BB'C > \angle BAC \geq \angle BCA$, so $BB' < BC = a$.

 Also, $\angle CC'B > \angle BAC \geq \angle ABC$, so $CC' < BC = a$.

 Thus $s = AA' + BB' + CC' < b + a + a = \boxed{\textbf{(A)}\ 2a + b}$.

 -edited by j314andrews

 



\textbf{Solution 2 (Answer Choices)}

This question has a unique correct answer, so if choice (X) is always less than or equal to choice (Y), then (X) must be incorrect, as otherwise (Y) would also be correct.

 Since $c \leq b$, $2a + c \leq 2a + b$, so (B) is incorrect.

 Since $b \leq a$, $2b + c \leq 2a + c$, so (C) is incorrect.

 Since $b \leq a$, $b + a + b \leq a + a + b$, that is, $a + 2b \leq 2a + b$.  So (D) is incorrect.

 Since $c \leq a$, $c + a + b \leq a + a + b$, that is, $a + b + c \leq 2a + b$.  So (E) is incorrect.

 Therefore, by process of elimination $\boxed{\textbf{(A)}\ 2a + b}$ is correct.

 -j314andrews

 


	\item (\textit{AHSME 1964}) A watch loses $2\frac{1}{2}$ minutes per day. It is set right at $1$ P.M. on March $15$. 
Let $n$ be the positive correction, in minutes, to be added to the time shown by the watch at a given time. 
When the watch shows $9$ A.M. on March $21$, $n$ equals:


 $\textbf{(A) }14\frac{14}{23}\qquad \textbf{(B) }14\frac{1}{14}\qquad \textbf{(C) }13\frac{101}{115}\qquad \textbf{(D) }13\frac{83}{115}\qquad \textbf{(E) }13\frac{13}{23}$

 \textbf{Answer: }A



\textbf{Solution}

From March 15 $1$ P.M. on the watch to March 21 $9$ A.M. on the watch, the watch passed $20 + 5 \times 24 = 140$ hours. 

 Since $1$ watch hour equals $\frac{24}{23 + \frac{57.5}{60}} = \frac{576}{575}$ real hour, the difference between the watch time and the actual time passed is $140 \times \left( \frac{576}{575} - 1 \right) = \frac{28}{115}$ hour $=14\frac{14}{23}$ minutes.

 


\end{enumerate}